\chapter{System Design and Implementation}\label{ch:system}

This chapter describes how the methodology of \cref{ch:method} was turned into a working system. It starts from
the requirements, presents the layered architecture and the design of its modules, the dimensional data store
and the interactive dashboard, and then explains how the system is deployed, tested and reproduced.

\section{Requirements analysis}
The requirements were derived from the objectives (O1--O7) and the evaluation criteria of the project proposal.
Following \citet{sommerville2016}, they are divided into functional requirements, which state what the system
must do, and non-functional requirements, which state the qualities it must have
(\cref{tab:funreq,tab:nonfunreq}). Each requirement is traced to the part of the system that satisfies it and to
the evidence that it was met.

\begin{table}[htbp]\centering\small
\caption{Functional requirements and their implementation}\label{tab:funreq}
\begin{tabularx}{\linewidth}{@{}lX>{\raggedright\arraybackslash}p{3.7cm}@{}}
\toprule
\textbf{ID} & \textbf{Requirement} & \textbf{Implemented in} \\ \midrule
FR1 & Load the raw transaction file with explicit types and validate it against schema, domain and identity rules & \texttt{data.py} \\
FR2 & Repair missing, duplicate and impossible values and flag statistical outliers, with a cleaning log & \texttt{data.py} \\
FR3 & Store cleaned data in a queryable dimensional model & \texttt{data.py}, SQLite \\
FR4 & Produce descriptive statistics, trend, seasonality, concentration and hypothesis tests & \texttt{eda.py} \\
FR5 & Build leakage-safe features for monthly product demand & \texttt{features.py} \\
FR6 & Train, tune and compare regression, classification, clustering and anomaly models against benchmarks & \texttt{forecasting.py}, \texttt{classification.py}, \texttt{clustering.py}, \texttt{anomaly.py} \\
FR7 & Forecast twelve months ahead with prediction intervals & \texttt{forecasting.py} \\
FR8 & Generate plain-language insights with recommended actions & \texttt{insights.py} \\
FR9 & Present KPIs, analyses and forecasts in an interactive, filterable web dashboard & \texttt{dashboard/app.py} \\
FR10 & Export every figure, table and number used by the report and website & \texttt{export.py}, \texttt{website.py} \\ \bottomrule
\end{tabularx}
\end{table}

\begin{table}[htbp]\centering\small
\caption{Non-functional requirements and how they were verified}\label{tab:nonfunreq}
\begin{tabularx}{\linewidth}{@{}lp{3.1cm}X@{}}
\toprule
\textbf{ID} & \textbf{Quality} & \textbf{Target and verification} \\ \midrule
NFR1 & Performance & Validate and clean \NRows rows in under one minute on a laptop; measured under 20 s \\
NFR2 & Reproducibility & One command regenerates all outputs; fixed random seed; verified by repeated runs \\
NFR3 & Correctness & Warehouse totals reconcile with the cleaned data to 0.00 BDT; nine automated tests pass \\
NFR4 & Usability & Runs in any browser without installation; global filters; plain-language captions \\
NFR5 & Availability and cost & Public HTTPS hosting on free platforms; automatic redeployment on every push \\
NFR6 & Maintainability & Modules with a single responsibility; one configuration file; documented code \\
NFR7 & Privacy & No personal data; raw file excluded from the public repository \\ \bottomrule
\end{tabularx}
\end{table}

\begin{figure}[htbp]\centering
\begin{tikzpicture}[every node/.style={font=\footnotesize},
  actor/.style={draw, rounded corners, fill=gray!15, minimum width=2.5cm, minimum height=9mm, align=center},
  uc/.style={draw, ellipse, minimum width=3.6cm, minimum height=8mm, align=center, fill=white},
  arr/.style={thick}]
\draw[rounded corners, dashed] (2.85,-4.55) rectangle (10.15,3.8);
\node[font=\footnotesize\bfseries] at (6.5,3.4) {Dan Cake analytics framework};
\foreach \i/\t in {1/Run pipeline,2/Validate and clean data,3/Train and evaluate models,4/Regenerate report and website,
                   5/Explore KPIs and trends,6/View demand forecast,7/Check anomalies,8/Read insights and actions,9/Download forecast CSV}
  \node[uc, minimum width=4.6cm] (u\i) at (6.5,{2.75-0.85*(\i-1)}) {\t};
\node[actor] (an) at (0.6,1.475) {Data analyst};
\node[actor] (mg) at (12.4,-2.2) {Manager};
\foreach \i in {1,...,4} \draw[arr] (an.east) -- (u\i.west);
\foreach \i in {5,...,9} \draw[arr] (mg.west) -- (u\i.east);
\end{tikzpicture}
\caption{Use-case view of the framework: the analyst operates the pipeline, managers consume its results}\label{fig:usecase}
\end{figure}

Two groups of users interact with the system (\cref{fig:usecase}). The \emph{data analyst} runs the pipeline,
inspects the validation report and maintains the models; \emph{managers} use the dashboard and the website to
monitor performance, read forecasts and act on the recommendations.

\section{Architecture of the framework}
The framework is organised in layers that mirror the workflow in the proposal (\cref{fig:architecture}). Each
layer is a Python module with a single responsibility. The modules communicate only through files with stable
formats (Parquet, SQLite and JSON), so any layer can be replaced without changing the others. For example, the
Streamlit dashboard could be swapped for Power BI reading the same SQLite store.

\begin{figure}[htbp]\centering
\begin{tikzpicture}[node distance=3.5mm, every node/.style={font=\footnotesize},
  L/.style={draw=brand, rounded corners, fill=brand!7, minimum width=12.5cm, minimum height=8.5mm, align=left, text width=12cm},
  arr/.style={-{Stealth[length=2mm]}, thick, brand}]
\node[L] (a) {\textbf{1. Data sources} \hfill raw CSV export (\NRows transactions, \NCols columns)};
\node[L, below=of a] (b) {\textbf{2. Ingestion \& validation} \hfill typed load, schema, rule and identity checks \hfill \texttt{data.py}};
\node[L, below=of b] (c) {\textbf{3. Cleaning \& storage} \hfill imputation, de-duplication, outlier flags, star schema (SQLite) \hfill \texttt{data.py}};
\node[L, below=of c] (d) {\textbf{4. Exploratory analytics} \hfill trends, seasonality, tests, elasticity \hfill \texttt{eda.py}};
\node[L, below=of d] (e) {\textbf{5. Feature engineering} \hfill leakage-safe lags, rolling statistics, calendar, price \hfill \texttt{features.py}};
\node[L, below=of e] (f) {\textbf{6. Machine learning} \hfill forecasting $\cdot$ classification $\cdot$ clustering $\cdot$ anomaly detection};
\node[L, below=of f] (g) {\textbf{7. Evaluation \& explanation} \hfill benchmarks, metrics, permutation importance, SHAP};
\node[L, below=of g] (h) {\textbf{8. Presentation} \hfill Streamlit dashboard $\cdot$ project website $\cdot$ \LaTeX{} report \hfill \texttt{dashboard/}};
\node[L, below=of h, fill=accent!10, draw=accent] (i) {\textbf{9. Decision support} \hfill automatically generated insights and recommendations \hfill \texttt{insights.py}};
\foreach \x/\y in {a/b,b/c,c/d,d/e,e/f,f/g,g/h,h/i} \draw[arr] (\x)--(\y);
\end{tikzpicture}
\caption{Layered architecture of the intelligent analytics framework}\label{fig:architecture}
\end{figure}

A single command, \texttt{python run\_pipeline.py}, executes all layers in order. It writes:
\begin{itemize}
  \item the cleaned dataset (\path{data/processed}) and the warehouse (\path{data/warehouse/dancake.db});
  \item metrics in JSON (\path{outputs/metrics}) and tables in CSV (\path{outputs/tables});
  \item every figure and \LaTeX{} table of this report (\path{report/figures}, \path{report/tables}), and a
        file of \LaTeX{} macros holding every number quoted in the text (\path{report/generated/numbers.tex});
  \item compact aggregated data for the dashboard and the website (\path{data/dashboard}).
\end{itemize}
Because no result is typed by hand, the report, dashboard and website cannot drift apart. A new month of data
can be incorporated by re-running one command.

\section{Module design}
Following the principle of high cohesion and low coupling \citep{pressman2020}, each module has one
responsibility and exposes a small number of functions (\cref{tab:modules}). Shared settings (file paths, column
names, the train/test boundary and the random seed) live in a single configuration module. This avoids the
``configuration debt'' that \citet{sculley2015} identify as a common weakness of machine-learning systems. The
whole framework, including the dashboard, comprises about 3,900 lines of Python.

\begin{table}[htbp]\centering\small
\caption{Modules of the framework (package \texttt{src/dancake})}\label{tab:modules}
\begin{tabularx}{\linewidth}{@{}lrX@{}}
\toprule
\textbf{Module} & \textbf{Lines} & \textbf{Responsibility} \\ \midrule
\texttt{config.py} & 50 & paths, column names, train/test dates, random seed \\
\texttt{data.py} & 220 & typed loading, validation report, cleaning, outlier flags, star-schema load \\
\texttt{eda.py} & 304 & KPIs, trends, seasonality, Pareto/ABC, hypothesis tests, elasticity, EDA figures \\
\texttt{features.py} & 86 & leakage-safe monthly panel features (lags, rolling, calendar, price, product) \\
\texttt{forecasting.py} & 361 & seven regressors, naive and seasonal-naive baselines, Holt--Winters, SARIMA, recursive test, 12-month outlook \\
\texttt{classification.py} & 191 & seven classifiers for large orders, threshold choice, permutation importance \\
\texttt{clustering.py} & 203 & product and market segmentation, validity indices, cluster naming \\
\texttt{anomaly.py} & 138 & daily residual monitoring and Isolation Forest transaction screening \\
\texttt{insights.py} & 85 & rule-based generation of insights and recommended actions \\
\texttt{export.py} & 156 & \LaTeX{} tables, number macros and dashboard data \\
\texttt{website.py} & 166 & static project website \\
\texttt{utils.py} & 89 & shared helpers: figure saving, JSON I/O, currency formatting, \LaTeX{} table writer \\
\texttt{run\_pipeline.py} & 46 & orchestration: runs all layers in order \\
\texttt{dashboard/app.py} & 383 & Streamlit dashboard with nine pages \\ \bottomrule
\end{tabularx}
\end{table}

\Cref{fig:sequence} shows the order in which the orchestrator calls the modules and the artefacts that pass
between them.

\begin{figure}[htbp]\centering
\begin{tikzpicture}[xscale=0.96, every node/.style={font=\footnotesize},
  life/.style={draw, fill=gray!12, minimum width=2.25cm, minimum height=8mm, align=center},
  msg/.style={-{Stealth[length=1.8mm]}, thick}]
\foreach \x/\n/\t in {0/r/run\_pipeline,2.55/d/data,5.1/e/eda,7.65/f/features+ML,10.2/x/export,12.75/w/dashboard}{
  \node[life] (\n) at (\x,0) {\texttt{\t}};
  \draw[dashed, gray] (\x,-0.45) -- (\x,-6.5);}
\draw[msg] (0,-0.95) -- node[above]{load + validate} (2.55,-0.95);
\draw[msg] (2.55,-1.65) -- node[above]{clean.parquet, warehouse.db} (5.1,-1.65);
\draw[msg] (0,-2.35) -- node[above, pos=0.7]{explore} (5.1,-2.35);
\draw[msg] (5.1,-3.05) -- node[above]{eda.json, figures} (10.2,-3.05);
\draw[msg] (0,-3.75) -- node[above, pos=0.8]{features, train, evaluate} (7.65,-3.75);
\draw[msg] (7.65,-4.45) -- node[above]{metrics.json, forecasts} (10.2,-4.45);
\draw[msg] (0,-5.15) -- node[above, pos=0.85]{export} (10.2,-5.15);
\draw[msg] (10.2,-5.85) -- node[above, align=center]{tables, numbers,\\ dashboard data} (12.75,-5.85);
\end{tikzpicture}
\caption{Sequence of calls in one run of the pipeline and the artefacts exchanged}\label{fig:sequence}
\end{figure}

\section{Dimensional data store}
The cleaned data are loaded into a star schema in SQLite (\cref{fig:star}). The fact table \texttt{fact\_sales}
holds one row per transaction with its measures. The dimensions describe products, dates (including a
Friday--Saturday weekend flag, as in Bangladesh) and the 220 market segments formed by region, channel and
customer type. Indexes on the foreign keys make typical KPI queries return in well under a second. A
reconciliation check confirms that total net sales in the warehouse equal the total in the cleaned dataset to
the last paisa.

\begin{figure}[htbp]\centering
\begin{tikzpicture}[every node/.style={font=\footnotesize},
  T/.style={draw=brand, rectangle split, rectangle split parts=2, rounded corners, fill=white, text width=3.7cm, align=left},
  F/.style={T, draw=accent, fill=accent!6, text width=4.0cm}]
\node[F] (f) at (0,0) {\textbf{fact\_sales} (\NRows)\nodepart{second}transaction\_id\\date\_id (FK)\\product\_id (FK)\\segment\_id (FK)\\list/market/unit price\\discount\_pct, quantity\\gross, discount, net sales\\flag\_outlier\_qty};
\node[T] (p) at (-5.2,1.6) {\textbf{dim\_product} (\NProducts)\nodepart{second}product\_id (PK)\\product\_name, category\\pack\_size, launch\_year\\mrp\_bdt};
\node[T] (d) at (5.2,1.6) {\textbf{dim\_date} (2,465)\nodepart{second}date\_id (PK), date\\year, quarter, month\\day\_of\_week, is\_weekend};
\node[T] (s) at (0,-4.1) {\textbf{dim\_segment} (220)\nodepart{second}segment\_id (PK)\\region, channel\\customer\_type};
\draw[thick, brand] (f.west) -- (p.east); \draw[thick, brand] (f.east) -- (d.west); \draw[thick, brand] (f.south) -- (s.north);
\end{tikzpicture}
\caption{Star schema of the analytical data store}\label{fig:star}
\end{figure}

The star schema can be queried directly in SQL. \Cref{app:sql} lists analytical queries that reproduce the key figures of \cref{ch:eda}; they can also serve as the data source of a Power BI report.

\section{Interactive dashboard}\label{ch:dashboard}

\subsection{Design goals}
The dashboard is where the framework meets its users, the sales, supply-chain and senior managers. Following
\citet{few2013}, \citet{tufte2001} and \citet{knaflic2015}, it was designed to meet five requirements:
\begin{enumerate}
  \item \textbf{At-a-glance KPIs}: net sales, units, orders, average order value, discount rate and
        year-to-date growth appear at the top of the overview page;
  \item \textbf{Consistent filtering}: one set of filters (years, region, channel, customer type, category)
        applies across all pages;
  \item \textbf{The three management questions}: pages progress from \emph{what happened} (overview, trends,
        products, markets) to \emph{why} (drivers, segments) and \emph{what next} (forecast, anomalies,
        recommendations);
  \item \textbf{Honest presentation}: forecasts are always shown with prediction intervals, model accuracy is
        visible to users, and colour is used only to encode meaning. Following \citet{tufte2001}, decoration
        that carries no data (3-D effects, gauges, heavy grid lines) is avoided;
  \item \textbf{Free and accessible}: the dashboard runs in any web browser and is hosted at no cost.
\end{enumerate}

\subsection{Implementation}
The dashboard was built with Streamlit and Plotly in a single Python file (\path{dashboard/app.py}). It does
not read the 500,000 raw rows. It reads compact aggregates produced by the pipeline: a monthly cube by region
$\times$ channel $\times$ customer type $\times$ category, monthly product-by-region totals, forecasts, model
scores, segment tables and anomaly lists. Their total size is about 2~MB, so the application starts in seconds
on the free hosting tier, while every filter combination remains exact. Data are cached in memory after the
first load. Each page has its own web address (for example \path{?page=forecast}), so a manager can bookmark
or share a specific view. \Cref{tab:pages} lists the nine pages, and
\cref{fig:dashoverview,fig:dashtrends,fig:dashproducts,fig:dashmarkets,fig:dashforecast,fig:dashmodels,fig:dashanomalies,fig:dashinsights,fig:dashquality}
show each of them as deployed. All screenshots were captured automatically from the running application by
\path{tools/capture_screenshots.py}, so they can be refreshed whenever the data change.

\begin{table}[htbp]\centering\small
\caption{Pages of the interactive dashboard}\label{tab:pages}
\begin{tabularx}{\linewidth}{@{}lX@{}}
\toprule
\textbf{Page} & \textbf{Content} \\ \midrule
Executive overview & KPI cards, monthly sales with moving average, category, region and channel mix, annual table \\
Trends \& seasonality & Stacked trends by any dimension, year $\times$ month heat map, seasonal index, year-on-year lines \\
Products \& portfolio & Top-$N$ products, Pareto/ABC analysis, ML product segments, single-product drill-down \\
Markets \& customers & Region $\times$ channel heat maps, customer types, regional CAGR, ML market segments \\
Demand forecast & 12-month outlook with interval, category and product forecasts, model leaderboard, CSV download \\
Machine-learning models & Regression, classification and clustering results with feature importance \\
Anomaly monitor & Daily sales against expected level with flagged days; most suspicious orders \\
Insights \& recommendations & Automatically generated insights with recommended actions \\
Data quality \& method & Validation report, cleaning log and architecture diagram \\ \bottomrule
\end{tabularx}
\end{table}

\subsection{Walk-through of the dashboard}
\paragraph{What happened.} The \emph{executive overview} (\cref{fig:dashoverview}) opens with six KPI cards and
the monthly sales trend with its 12-month average. Below them, the category, region and channel mix and an
annual summary table answer the first questions a manager asks. The \emph{trends} page
(\cref{fig:dashtrends}) shows how any dimension develops over time and makes the December peak and February
trough visible in a year~$\times$~month heat map. The \emph{products} page (\cref{fig:dashproducts}) ranks the
products, draws the Pareto curve with ABC classes and lets the user drill down into a single product. The
\emph{markets} page (\cref{fig:dashmarkets}) compares regions, channels and customer types and shows the
market segments found by clustering.

\paragraph{What next.} The \emph{demand forecast} page (\cref{fig:dashforecast}) presents the 12-month outlook
with its 80\% prediction interval, forecasts by category and product, and the leaderboard of forecasting
models. The forecasts can be downloaded as CSV for the production plan. The \emph{machine-learning} page
(\cref{fig:dashmodels}) exposes the evidence behind the forecasts: test errors of all models against the
baselines, the recursive multi-step test and feature importance, plus the classification and clustering
results in separate tabs.

\paragraph{Control and action.} The \emph{anomaly monitor} (\cref{fig:dashanomalies}) plots daily sales against
the expected level, marks the flagged days and lists the most suspicious orders for checking. The
\emph{insights} page (\cref{fig:dashinsights}) turns all results into plain-language findings, each followed by
a recommended action. Finally, the \emph{data quality} page (\cref{fig:dashquality}) documents the validation
checks and the cleaning log, so users can see why the numbers can be trusted.

% One full-width screenshot per dashboard page (the upper part of each page, captured at double
% resolution by tools/capture_screenshots.py) so that the text in the screenshots remains readable.

\begin{figure}[htbp]\centering
\includegraphics[width=\linewidth,height=.42\textheight,keepaspectratio]{screens/dash_overview_report}
\caption{Dashboard page 1 -- Executive overview: KPI cards and monthly sales trend}\label{fig:dashoverview}
\end{figure}

\begin{figure}[htbp]\centering
\includegraphics[width=\linewidth,height=.42\textheight,keepaspectratio]{screens/dash_trends_report}
\caption{Dashboard page 2 -- Trends and seasonality}\label{fig:dashtrends}
\end{figure}

\begin{figure}[htbp]\centering
\includegraphics[width=\linewidth,height=.42\textheight,keepaspectratio]{screens/dash_products_report}
\caption{Dashboard page 3 -- Products and Pareto/ABC analysis}\label{fig:dashproducts}
\end{figure}

\begin{figure}[htbp]\centering
\includegraphics[width=\linewidth,height=.42\textheight,keepaspectratio]{screens/dash_markets_report}
\caption{Dashboard page 4 -- Markets: region $\times$ channel heat maps}\label{fig:dashmarkets}
\end{figure}

\begin{figure}[htbp]\centering
\includegraphics[width=\linewidth,height=.42\textheight,keepaspectratio]{screens/dash_forecast_report}
\caption{Dashboard page 5 -- 12-month demand forecast with its 80\% interval}\label{fig:dashforecast}
\end{figure}

\begin{figure}[htbp]\centering
\includegraphics[width=\linewidth,height=.42\textheight,keepaspectratio]{screens/dash_models_report}
\caption{Dashboard page 6 -- Model errors against the baselines}\label{fig:dashmodels}
\end{figure}

\begin{figure}[htbp]\centering
\includegraphics[width=\linewidth,height=.42\textheight,keepaspectratio]{screens/dash_anomalies_report}
\caption{Dashboard page 7 -- Anomaly monitor}\label{fig:dashanomalies}
\end{figure}

\begin{figure}[htbp]\centering
\includegraphics[width=\linewidth,height=.42\textheight,keepaspectratio]{screens/dash_insights_report}
\caption{Dashboard page 8 -- Generated insights and recommendations}\label{fig:dashinsights}
\end{figure}

\begin{figure}[htbp]\centering
\includegraphics[width=\linewidth,height=.42\textheight,keepaspectratio]{screens/dash_quality_report}
\caption{Dashboard page 9 -- Data quality and method}\label{fig:dashquality}
\end{figure}

\begin{figure}[htbp]\centering
\includegraphics[width=\linewidth,height=.42\textheight,keepaspectratio]{screens/website}
\caption{Project website on GitHub Pages, linking the dashboard, report and code}\label{fig:website}
\end{figure}

\section{Deployment and operation}\label{sec:deployment}
\subsection{Live system}
The complete project is publicly available on three free platforms (\cref{tab:links}). The source code is
version-controlled on GitHub. The interactive dashboard runs on Streamlit Community Cloud, and the project
website, which also serves this report as a PDF, is hosted on GitHub Pages. All three use HTTPS. The raw
transaction file is deliberately excluded from the public repository by \texttt{.gitignore}. Only the small,
aggregated files that the dashboard needs are published, and commits use a private GitHub e-mail address.

\begin{table}[htbp]\centering\small
\caption{Public addresses of the deployed system}\label{tab:links}
\begin{tabularx}{\linewidth}{@{}lXl@{}}
\toprule
\textbf{Component} & \textbf{Address} & \textbf{Platform} \\ \midrule
Interactive dashboard & \dashurl & Streamlit Community Cloud \\
Project website and PDF report & \siteurl & GitHub Pages \\
Source code and history & \repo & GitHub \\ \bottomrule
\end{tabularx}
\end{table}

\subsection{Continuous deployment}
Deployment is automatic (\cref{fig:cicd}). When a change is pushed to the \texttt{main} branch, Streamlit Cloud
rebuilds the dashboard from \path{dashboard/app.py} and its small \path{dashboard/requirements.txt}, and GitHub
Pages republishes the \path{docs/} folder. A new month of data therefore reaches managers through three steps:
re-run the pipeline, commit, and push. No server needs to be maintained. On the free tier, the dashboard sleeps
after several days without visitors and wakes in about half a minute on the next visit, which is acceptable for
a monthly planning tool.

\begin{figure}[htbp]\centering
\begin{tikzpicture}[every node/.style={font=\footnotesize}, scale=0.92, transform shape,
  box/.style={draw=brand, rounded corners, fill=brand!8, minimum width=2.45cm, minimum height=10mm, align=center},
  web/.style={box, draw=accent, fill=accent!10},
  arr/.style={-{Stealth[length=2mm]}, thick, brand}]
\node[box] (data) at (0,0) {Raw sales data\\(local, private)};
\node[box] (pipe) at (3.2,0) {\texttt{run\_pipeline.py}\\models, figures, report};
\node[box] (git)  at (7.35,0) {GitHub repository\\\texttt{main} branch};
\node[web] (st)   at (10.55,0.9) {Streamlit Cloud\\dashboard};
\node[web] (gp)   at (10.55,-0.9) {GitHub Pages\\website + PDF};
\node[box, minimum width=1.9cm] (users) at (13.45,0) {Managers\\(browser)};
\draw[arr] (data) -- (pipe);
\draw[arr] (pipe) -- node[above=0.5mm, align=center, font=\scriptsize]{git\\push} (git);
\draw[arr] (git.east) -- ++(0.3,0) |- (st.west);  \draw[arr] (git.east) -- ++(0.3,0) |- (gp.west);
\draw[arr] (st.east) -- ++(0.35,0) |- (users.west); \draw[arr] (gp.east) -- ++(0.35,0) |- (users.west);
\end{tikzpicture}
\caption{Continuous deployment: one push updates both the dashboard and the website}\label{fig:cicd}
\end{figure}

\begin{figure}[htbp]\centering
\includegraphics[width=.85\linewidth]{screens/github_repo}
\caption{Public GitHub repository holding the code, the dashboard data, the website and the report}\label{fig:github}
\end{figure}

If a custom domain name is bought later, it can be attached to GitHub Pages through a DNS \texttt{CNAME} record
without any change to the code. \Cref{app:reproduce} gives the step-by-step deployment procedure.

\section{Testing and quality assurance}\label{sec:testing}
Testing followed the practice recommended by \citet{sommerville2016}: automated unit tests for the components
whose failure would silently corrupt results, complemented by integration checks inside the pipeline and by
acceptance checks against the requirements.

\paragraph{Unit tests.} The test suite (\path{tests/test_pipeline.py}, run with \texttt{pytest}) contains nine
test cases. Each builds a small synthetic sample, damages it in a known way and asserts that the pipeline
behaves correctly (\cref{tab:tests-unit}). All nine tests pass in about two seconds.

\begin{table}[htbp]\centering\small
\caption{Automated unit tests and the requirement each one protects}\label{tab:tests-unit}
\begin{tabularx}{\linewidth}{@{}>{\raggedright\arraybackslash\footnotesize\ttfamily}p{4.3cm}Xl@{}}
\toprule
\normalfont\small\textbf{Test} & \textbf{What it asserts} & \textbf{Req.} \\ \midrule
clean\_data\_\allowbreak passes\_validation & clean data show zero missing cells, zero duplicate IDs and exact accounting identities & FR1 \\
missing\_values\_\allowbreak are\_derived\_exactly & 60 deleted quantities are restored exactly from the identities & FR2 \\
duplicates\_and\_\allowbreak impossible\_values\_\allowbreak are\_fixed & exact and key duplicates are removed; a 250\% discount is repaired & FR2 \\
outliers\_are\_\allowbreak flagged\_not\_deleted & a 400-unit order is kept and flagged & FR2 \\
pack\_grams (4 cases) & pack sizes such as ``30g x 12'' are parsed to grams & FR5 \\
features\_do\_not\_\allowbreak use\_the\_current\_month & lag and rolling features use only past months; changing the current month does not change its features & FR5, NFR3 \\ \bottomrule
\end{tabularx}
\end{table}

\paragraph{Integration checks.} Each run of the pipeline checks its own outputs. The validation report is
regenerated; total net sales in the SQLite warehouse are reconciled with the cleaned dataset (difference
0.00 BDT); and every model is scored against its baseline, so a broken feature would show up immediately as a
model that no longer beats the naive rule.

\paragraph{Acceptance checks.} The deployed dashboard was opened page by page in a browser and every page was
captured automatically by \path{tools/capture_screenshots.py}; the screenshots in \cref{ch:dashboard} are the
result. The non-functional targets in \cref{tab:nonfunreq} were checked by timing the pipeline and by
regenerating the complete report from a clean state.

\section{Chapter summary}
The framework was implemented as a layered set of single-purpose Python modules, orchestrated by one command,
with a star-schema warehouse, a nine-page Streamlit dashboard and a generated website and report. Ten functional
and seven non-functional requirements were defined and traced to the implementation, and correctness is
protected by automated tests and reconciliation checks. The system is deployed publicly at no cost with
continuous deployment. \Cref{ch:results} now reports what the system found in the Dan Cake data and how well its
models perform.
